% ECONOMICS_PROBLEM: {"id": "L40-616", "title": "Effectiveness of antitrust remedies", "jel_code": "L40", "source_id": "OP-0080"}
% !TeX program = xelatex
\documentclass[11pt,a4paper]{article}
\usepackage[margin=23mm]{geometry}
\usepackage{fontspec}
\setmainfont{Latin Modern Roman}
\usepackage{amsmath,amssymb,mathtools}
\usepackage{xcolor,enumitem,booktabs,longtable,array,xurl,hyperref}
\definecolor{muted}{HTML}{53636C}
\hypersetup{colorlinks=true,urlcolor=blue,linkcolor=blue}
\setlength{\parindent}{0pt}
\setlength{\parskip}{5pt}
\newcommand{\R}{\mathbb R}
\newcommand{\E}{\mathbb E}
\newcommand{\N}{\mathbb N}
\newcommand{\ind}{\mathbf 1}
\newcommand{\Var}{\operatorname{Var}}
\newcommand{\Cov}{\operatorname{Cov}}
\newcommand{\supp}{\operatorname{supp}}
\DeclareMathOperator*{\argmin}{arg\,min}
\DeclareMathOperator*{\argmax}{arg\,max}
\newcommand{\entrymeta}[3]{{\sffamily\small\color{muted}\textbf{\texttt{#1}}\quad|\quad #2\par #3\par}}
\newcommand{\Problemref}[1]{\texttt{#1}}
\newcommand{\JELref}[2]{\texttt{#2} (JEL \texttt{#1})}
\begin{document}
\section*{Effectiveness of antitrust remedies}
\textbf{Problem ID:} \texttt{L40-616}\quad \textbf{JEL:} \texttt{L40}\par
% BEGIN ECONOMICS STATEMENT
\label{problem:OP-0080}
\label{atlas:prob:I10}

\noindent\textbf{Original identifier:} I10 (atlas item 80). \textbf{Classification:} Remedy effectiveness causal evaluation research program. \textbf{Original tractability label:} Medium (editorial, not a complexity result).

\subsection*{Definitions and assumptions}

Let a case $j$ have predecision covariates $X_j$ and eligible remedies $r\in\mathcal R_j$, including divestiture bundles, conduct restrictions, and interoperability obligations. The counterfactual $Y_{jt}(r)$ records quality-adjusted prices, output, entry, and service quality over a declared horizon $t=1,\ldots,T$. Remedy assignment is generally selected by the case, negotiations, and anticipated harm. Define enforcement resources, compliance monitoring, and buyer suitability as components of the regime. A target population of cases is specified independently of which cases were studied. Noncompliance is a post-remedy response, not a baseline confounder to control away automatically.

\subsection*{Formal research question}

Identify or bound policy contrasts $\tau_t(r,r')=\mathbb E_{j\sim F_*}[Y_{jt}(r)-Y_{jt}(r')]$ and a restoration criterion relative to an identified no-merger competitive baseline $Y_{jt}(0)$. For a prespecified discrepancy function $d$, assess
\[
\Pr_{j\sim F_*}\!\left(\max_{1\leq t\leq T}d(Y_{jt}(r),Y_{jt}(0))\leq\varepsilon\right),
\]
with a declared tolerance $\varepsilon$. Determine conditions for causal identification using randomized implementation where feasible, credible matched or synthetic comparison markets, or structural restrictions. Compare remedies by consumer welfare net of implementation and monitoring resource costs, with explicit selection sensitivity.

\subsection*{Interpretation and scope}

Durable restoration means success over a stated finite horizon; finite data cannot prove that firms will never circumvent a remedy. A remedy can improve outcomes compared with an unremedied merger while still failing to restore the no-merger baseline. These are different contrasts and should be reported separately. Divestitures require a viable buyer, retained capabilities, and access to complementary assets; counting assets transferred does not establish competitive effectiveness. Conduct remedies depend on monitoring, enforceability, and firms' adaptation, making implementation part of the treatment rather than incidental detail. Case studies selected after unusual transactions do not automatically generalize to ordinary or blocked cases. Synthetic controls require credible counterfactual fit and no contaminating spillovers, and selection on unobserved anticipated harm can defeat matched comparisons. A useful extension would register the remedy, target outcomes, comparison strategy, and data access before implementation, then follow entrants and consumers beyond the immediate transaction. Cross-case comparison must allow heterogeneous industries and remedy objectives rather than rank all remedies by one unadjusted price average.

\subsection*{Source audit and status}

Farrell (2003) is resolved to the verified Berkeley working-paper version [S1], avoiding an unverified book-chapter edition. [S2] is Kwoka's 2015 retrospective monograph. ``FTC merger-retrospective literature'' is not a complete citation; [S3] supplies one explicit 2017 staff study without claiming it uniquely exhausts that phrase. These references contain existing remedy evidence. The target population and restoration test are editorial extensions, with causal and generalization limits stated explicitly.

\subsection*{References}

\begin{enumerate}[label={[S\arabic*]},leftmargin=*]

\item Joseph Farrell (2003). \emph{Negotiation and Merger Remedies: Some Problems}. UC Berkeley Competition Policy Center working paper, August 2003. \url{https://escholarship.org/uc/item/9p72k8fn}\par \textit{Source role:} Economic analysis of remedy negotiations; working-paper version used. \textit{Locator:} Abstract and article metadata.

\item John Kwoka (2015). \emph{Mergers, Merger Control, and Remedies: A Retrospective Analysis of U.S. Policy}. MIT Press. \url{https://mitpress.mit.edu/9780262028486/mergers-merger-control-and-remedies/}\par \textit{Source role:} Synthesis of merger retrospectives. \textit{Locator:} Abstract and article metadata.

\item Federal Trade Commission, Bureaus of Competition and Economics (2017). \emph{The FTC's Merger Remedies 2006--2012: A Report of the Bureaus of Competition and Economics}. FTC staff report, January 2017. \url{https://www.ftc.gov/reports/ftcs-merger-remedies-2006-2012-report-bureaus-competition-economics}\par \textit{Source role:} Explicit replacement for unspecified FTC retrospective literature. \textit{Locator:} Report overview and case-study methodology.

\end{enumerate}

\subsection*{Common conventions: Source mathematical conventions}

Unless an entry states otherwise, agent and firm sets are finite. State and action spaces are measurable, strategies and outcome maps are measurable, probability laws are specified, and expectations exist for the targets under discussion. Infinite-horizon discount factors lie in $(0,1)$ and discounted objectives are finite. Norms, tolerances and complexity input sizes must be specified for computational comparisons. Constants or primitives that appear locally are inputs, not empirical facts asserted by this catalogue.

\textbf{Identification.} A statistical model $m\in\mathcal M$ implies an observable law $P_m$ and target $\theta(m)$. At law $P$, its identified set is
\[
\Theta(P;\mathcal M)=\{\theta(m):m\in\mathcal M,\ P_m=P\}.
\]
Point identification holds when this set is a singleton, or equivalently when $P_{m_1}=P_{m_2}$ implies $\theta(m_1)=\theta(m_2)$ for all compatible models. Sharp bounds mean bounds attained, or approached when the boundary is not attained, by compatible models. Writing this set is a definition, not a solution: the substantive task is to establish informative restrictions and derive or estimate the set. An unrestricted-confounding model may give only trivial bounds.

\textbf{Inference.} Identification, estimability and valid uncertainty quantification are separate questions. Fix $\alpha\in(0,1)$. A confidence procedure $C_n$ over a declared law class $\mathcal P$ has asymptotic uniform coverage at level $1-\alpha$ when
\[
\liminf_{n\to\infty}\inf_{P\in\mathcal P}\Pr_P\{\theta(P)\in C_n\}\geq1-\alpha.
\]
For set-identified targets the coverage event must be defined separately, for example coverage of the full identified set. If an entry uses identification-indexed models instead of uniquely indexed laws, the same coverage convention is applied over those models. Pointwise limits do not establish uniform validity, and asymptotic validity does not guarantee finite-sample precision. Sample sizes, dependence structures and weak-identification sequences are part of each application.

\textbf{Counterfactuals.} $Y_i(a)$ is a potential outcome under a specified intervention, including its timing, implementation and versions. It is not $Y_i$ conditional on voluntarily choosing $a$. A full intervention vector permits interference; shorthand scalar treatment notation does not silently impose no interference. Assignment, selection, attrition, anticipation and measurement must be specified. A claim about an instrument's local effect is not automatically a population policy effect.

\textbf{Economic equilibrium.} A counterfactual model includes best responses, constraints, feasibility, market clearing, state transitions and expectations consistent with the stated information. When several equilibria exist, either a declared selector is an input or the target is set-valued. Comparisons across policy regimes must use the stated selection convention. Reduced-form relationships are not presumed invariant after an equilibrium-changing intervention.

\textbf{Optimization and welfare.} Optimization is over the stated feasible set. If attainment is not established, $\sup$ or $\inf$ and $\varepsilon$-optimal choices replace an asserted maximizer or minimizer. For example, with value $V=\sup_{a\in\mathcal A}W(a)<\infty$, an $\varepsilon$-optimal set is $\{a\in\mathcal A:W(a)\geq V-\varepsilon\}$ for $\varepsilon>0$. Benefits, costs, risk and health or environmental outcomes can be aggregated only using declared units and conversion weights. Interpersonal utility weights, the treatment of fiscal transfers and foreign profits, discounting, and the behavioral welfare criterion are explicit inputs. A comparative-static derivative additionally requires existence and appropriate differentiability of the selected counterfactual.

\textbf{Sources.} Local labels [S1], [S2], and so forth restart in every question. Locators identify the relevant abstract, model, section or result at the level actually checked; a bibliographic or abstract-level verification is not represented as inspection of an entire proof. Working papers are distinguished from later publications and changing versions. Source summaries are paraphrased. All URL checks and editorial formulations refer to the audit date stated above.
% END ECONOMICS STATEMENT
\end{document}
